% ================================================================ TITLE PAGE
\begin{titlepage}
\newgeometry{left=1in,right=1in,top=0.9in,bottom=0.9in}
\centering
\setstretch{1.15}
{\Large\bfseries B.TECH.\ PROJECT ON\par}
\vspace{0.7cm}
{\huge\bfseries Reinforcement Learning for Swarm\\[4pt] Technologies of Drone\par}
\vspace{1.0cm}
{\large\itshape Submitted By\par}
\vspace{0.3cm}
{\large
\begin{tabular}{l@{\hspace{1.2cm}}l}
Nishant Raj    & 2023UEC2606\\
Sumit Kumar    & 2023UEC2554\\
Dheeraj Sharma & 2023UEC2599\\
Tarun Kumar    & 2023UEC2551\\
\end{tabular}\par}
\vspace{0.9cm}
{\large\itshape Under the Guidance of\par}
\vspace{0.25cm}
{\Large\bfseries \Supervisor\par}
{\normalsize \SupervisorDesignation, \Department\par}
\vfill
\includegraphics[width=0.19\textwidth]{nsut_logo}\par
\vfill
{\normalsize Project-I in partial fulfillment of requirement for the award of\par}
{\large\bfseries \Programme\par}
\vspace{0.8cm}
{\large\bfseries \Department\par}
{\large\bfseries NETAJI SUBHAS UNIVERSITY OF TECHNOLOGY\par}
{\large\bfseries NEW DELHI -- 110078\par}
\vspace{0.5cm}
{\large \SubmissionMonth\par}
\restoregeometry
\end{titlepage}

\pagenumbering{roman}
\setcounter{page}{1}

% ================================================================ CERTIFICATE
\chapter*{CERTIFICATE}
\addcontentsline{toc}{chapter}{Certificate}
\begingroup
\setstretch{1.6}
\setlength{\parskip}{12pt}
\sloppy
\vspace*{0.6cm}

Certified that \textbf{Nishant Raj (2023UEC2606)}, \textbf{Sumit Kumar (2023UEC2554)}, \textbf{Dheeraj Sharma
(2023UEC2599)} and \textbf{Tarun Kumar (2023UEC2551)} have carried out their project work presented in this
project entitled \textbf{``\ReportTitle''} for the award of Bachelor of Technology, \Department, \University,
New Delhi, under my supervision.

The project embodies results of original work, and studies are carried out by the students themselves, and the
contents of the project do not form the basis for the award of any other degree to the candidates or to anybody
else from this or any other University/Institution.

\vfill
\noindent
\begin{minipage}[b]{0.45\textwidth}
Place: New Delhi\\[10pt]
Date: \rule{3.5cm}{0.4pt}
\end{minipage}\hfill
\begin{minipage}[b]{0.48\textwidth}
\centering
\rule{5.5cm}{0.4pt}\\[4pt]
\textbf{\Supervisor}\\
\SupervisorDesignation\\
\Department\\
\University, New Delhi
\end{minipage}
\vspace*{1.2cm}
\endgroup

% ================================================================ DECLARATION
\chapter*{CANDIDATES' DECLARATION}
\addcontentsline{toc}{chapter}{Candidates' Declaration}
\begingroup
\setstretch{1.6}
\setlength{\parskip}{12pt}
\vspace*{0.6cm}

We hereby declare that the work presented in this project report entitled \textbf{``\ReportTitle''}, submitted
in partial fulfillment of the requirement for the award of the degree of Bachelor of Technology in Electronics
and Communication Engineering, \University, New Delhi, is an authentic record of our own work carried out under
the guidance of \textbf{\Supervisor}.

The matter embodied in this report has not been submitted by us for the award of any other degree. Wherever the
work of others has been used, it has been duly acknowledged and cited.

\vfill
\noindent
\begin{tabular*}{\textwidth}{@{\extracolsep{\fill}}cc}
\rule{5cm}{0.4pt} & \rule{5cm}{0.4pt}\\
\textbf{Nishant Raj} & \textbf{Sumit Kumar}\\
2023UEC2606 & 2023UEC2554\\[1.6cm]
\rule{5cm}{0.4pt} & \rule{5cm}{0.4pt}\\
\textbf{Dheeraj Sharma} & \textbf{Tarun Kumar}\\
2023UEC2599 & 2023UEC2551\\
\end{tabular*}

\vspace{1cm}
\noindent Place: New Delhi \hfill Date: \rule{3.5cm}{0.4pt}
\vspace*{1.2cm}
\endgroup

% ================================================================ ACKNOWLEDGEMENT
\chapter*{ACKNOWLEDGEMENT}
\addcontentsline{toc}{chapter}{Acknowledgement}
\begingroup
\setstretch{1.6}
\setlength{\parskip}{12pt}
\vspace*{0.6cm}

We would like to express our gratitude and appreciation to all those who made it possible to complete this
project.

Special thanks to our project supervisor, \textbf{\Supervisor}, whose guidance, stimulating suggestions and
encouragement helped us throughout the design, implementation and evaluation of this work.

We also sincerely thank the \Department, \University, for providing the facilities and environment required
for the project, and our colleagues for their valuable support and assistance throughout the project.

\vfill
\noindent
\begin{tabular*}{\textwidth}{@{\extracolsep{\fill}}cc}
\rule{5cm}{0.4pt} & \rule{5cm}{0.4pt}\\
\textbf{Nishant Raj} & \textbf{Sumit Kumar}\\[1.6cm]
\rule{5cm}{0.4pt} & \rule{5cm}{0.4pt}\\
\textbf{Dheeraj Sharma} & \textbf{Tarun Kumar}\\
\end{tabular*}
\vspace*{1.2cm}
\endgroup

% ================================================================ ABSTRACT
\chapter*{ABSTRACT}
\addcontentsline{toc}{chapter}{Abstract}
This work presents a reinforcement-learning system for the cooperative control of a drone swarm. Eight drones
follow a moving reference target in a fixed formation; when a non-cooperative object crosses the flight path,
the drones leave their formation slots, avoid the object and one another, and then return to \emph{their own}
slots relative to the moving target. A single neural policy is shared by all drones and each drone executes it
using only local observations of its own state, its six nearest neighbours and up to six nearby objects. The
policy combines a self encoder with attention-based set encoders, so that it is independent of the order in
which neighbours are observed, and outputs a desired three-dimensional velocity to a stabilised flight model.

The environment is implemented as a PettingZoo multi-agent environment with continuous swept collision
detection, a bounded per-drone reward and an explicit formation-recovery criterion. Training uses a three-stage
curriculum: in each stage the policy is first initialised from demonstrations of a sampling-based
model-predictive expert planner (noise-injected demonstrations followed by DAgger) and is then fine-tuned with
proximal policy optimisation (PPO) and generalised advantage estimation, while the best checkpoint is selected on
separate validation seeds.

The final model was evaluated on 140 held-out episodes (seven scenarios $\times$ 20 unseen seeds). It completed
138 episodes (98.6\%) without failure, achieved 100\% success in six of the seven scenarios, kept the
undisturbed formation error below 5\,cm (0.023--0.045\,m) and returned every drone to its own slot after every
disturbance in the completed episodes. A browser-based flight laboratory visualises the learned behaviour
interactively.

\vspace{0.5cm}
\noindent\textbf{Keywords:} multi-agent reinforcement learning, drone swarm, formation control, collision
avoidance, proximal policy optimisation, attention, imitation learning.

% ================================================================ LISTS (each on its own page)
\clearpage
{\small\setstretch{0.97}\tableofcontents}
\ifwithlists
\clearpage
\phantomsection
\addcontentsline{toc}{chapter}{List of Figures}
{\setstretch{1.1}\listoffigures}
\clearpage
\phantomsection
\addcontentsline{toc}{chapter}{List of Tables}
{\setstretch{1.1}\listoftables}
\fi
